\section*{ОБОЗНАЧЕНИЯ И СОКРАЩЕНИЯ}

БД – база данных.

ИС – информационная система.

ИТ – информационные технологии.

ПО – программное обеспечение.

РП – рабочий проект.

ТЗ – техническое задание.

ТП – технический проект.

ПО – программное обеспечение.

СУБД – система управления базами данных.

API (Application Programming Interface) – программный интерфейс приложения, набор правил и инструментов для взаимодействия программных компонентов.

HTTP (HyperText Transfer Protocol) – протокол передачи гипертекста.

HTTPS (HyperText Transfer Protocol Secure) – защищенный протокол передачи гипертекста.

SQL (Structured Query Language) – язык структурированных запросов к базам данных.

REST (Representational State Transfer) — архитектурный стиль взаимодействия компонентов распределённого приложения через HTTP.

JWT (JSON Web Token) — формат токена, используемый для аутентификации и авторизации пользователей.

JSON (JavaScript Object Notation) — текстовый формат обмена данными, применяемый для передачи информации между клиентской и серверной частями системы.

CRUD (Create, Read, Update, Delete) — основные операции по работе с данными: создание, чтение, обновление и удаление.

PDF (Portable Document Format) — формат электронных документов, используемый для генерации договоров аренды.
